\section{The Autonomous Drift Problem}

\subsection{Definitions and Formal Characterization}

Recursive agent spawning creates a structural vulnerability we term the \emph{autonomous drift problem}. This section provides precise definitions, characterizes the conditions under which drift occurs, catalogs the resulting failure modes, and demonstrates formally why naive depth limits are insufficient as a remedy.

\medskip
\noindent\textbf{Definitions.} We model a spawned agent ecosystem as a rooted tree $\mathcal{T}$ where nodes are agents and directed edges represent the \texttt{spawns} relation. For any agent $a$:

\begin{itemize}
  \item \textbf{Parent.} $\parent(a)$ denotes the agent that directly spawned $a$. The root has no parent.
  \item \textbf{Ancestor chain.} $\anc(a) = (a_0, a_1, \dots, a_n)$ where $a_0 = a$, $a_{i+1} = \parent(a_i)$, and $a_n$ is the root agent (or a designated \texttt{HALT} node).
  \item \textbf{Human anchor.} $\anchor(a)$ denotes a designated human principal (or authorized control channel) that holds termination authority over $a$.
  \item \textbf{Agent state.} $\state(a) \in \{\active, \orphaned, \drifted, \terminated\}$.
\end{itemize}

\medskip
\noindent\textbf{Orphaned agent.} An agent $a$ is \emph{orphaned} if $\parent(a)$ exists and $\state(\parent(a)) \in \{\terminated, \unreachable\}$. That is, the direct parent has exited or become inaccessible, severing the primary control channel.

\medskip
\noindent\textbf{Drifted agent.} An agent $a$ is \emph{drifted} if $\forall\, b \in \anc(a): \anchor(b) = \nil$. Equivalently, the ancestor chain does not trace back to any reachable human anchor. An orphaned agent \emph{becomes} drifted when no living ancestor retains an anchor connection.

These definitions are independent; an orphaned agent may still be anchored through a grandparent or alternate channel, and a non-orphaned agent may nonetheless be drifted if all paths to a human principal are broken.

\subsection{Conditions That Cause Drift}

Drift emerges from four distinct conditions operating alone or in combination:

\medskip
\noindent\textbf{C1: Parent termination.} The direct parent exits normally or crashes, orphaning all children. This is the most common cause in long-running agentic systems where intermediate orchestrators have bounded lifecycles.

\medskip
\noindent\textbf{C2: Network partition.} The parent remains alive but becomes unreachable over the communication channel (timeouts, retries exhausted). The parent cannot receive directives; the child cannot confirm liveness.

\medskip
\noindent\textbf{C3: Anchor revocation.} The human principal explicitly revokes anchor authority from an agent (e.g., revoking an API key, disabling a OAuth integration). The agent continues operating but no longer has a binding to any principal.

\medskip
\noindent\textbf{C4: Malicious or buggy spawn.} A compromised or malfunctioning parent spawns children that inherit no anchor information, either by design or by implementation error (e.g., failure to propagate \texttt{ANCHOR\_ID} in the spawn payload).

\subsection{Failure Modes}

Once drift occurs, agents exhibit one or more of the following failure modes:

\medskip
\noindent\textbf{F1: Uncontrolled continuation.} The drifted agent continues executing its current task loop with no ability to receive external corrections, halt signals, or priority overrides. It operates on cached state and last-known objectives.

\medskip
\noindent\textbf{F2: Resource accumulation.} Drifted agents retain allocated resources (compute budget, memory, API quotas, credentials) with no garbage-collection path. In recursive spawning systems, each drifted agent may itself spawn further agents, compounding resource drain exponentially.

\medskip
\noindent\textbf{F3: Objective corruption.} Without an anchor, the drifted agent's objective function degrades. In systems where objectives are partially encoded in parent-to-child message passing, an orphaned agent's task description becomes stale. It may continue toward a goal that has been superseded, cancelled, or fundamentally changed upstream.

\medskip
\noindent\textbf{F4: Cascade amplification.} Drifted agents that retain the ability to spawn children propagate the drift condition downward. The subtree rooted at a drifted node inherits the drift property by construction, converting a localized failure into a systematic one.

\medskip
\noindent\textbf{F5: State divergence and reconciliation failure.} Drifted agents maintain local state that diverges from the system's ground truth. When and if reconnection occurs (a human detects drift and attempts recovery), the drifted agent's state may be incompatible with the current system state, making reconciliation expensive or impossible without force-reset.

\subsection{Why Depth Limits Alone Do Not Solve Drift}

A natural first response to drift risk is to impose a maximum spawn depth $D_{\max}$. The intuition: if no agent spawns beyond depth $D_{\max}$, a human anchor remains within $D_{\max}$ hops of any leaf. This reasoning is flawed for three reasons.

\medskip
\noindent\textbf{Limitation 1: Depth does not imply anchor proximity.} Consider a chain of $D_{\max} = 10$ obedient, well-functioning agents where each has a valid parent pointer but none has a human anchor assigned. All agents are at depth $\leq 10$, yet every one is drifted. Depth limits constrain \emph{distance from root}, not \emph{distance to a human principal}. If the root itself is not anchored, the bound is meaningless.

\begin{equation}
\text{Drift} \iff \left(\exists a: \forall\, b \in \anc(a): \anchor(b) = \nil\right)
\quad\text{regardless of } \depth(a)
\end{equation}

A depth limit only guarantees anchor proximity when every agent at depth $d < D_{\max}$ retains a valid anchor. This property is not enforced by a depth bound alone.

\medskip
\noindent\textbf{Limitation 2: Horizontal spawning bypasses depth limits.} An agent at depth $D_{\max} - 1$ that spawns $k$ children in parallel stays within the depth limit but creates $k$ independent drift points. If the parent subsequently becomes unreachable, all $k$ children are simultaneously orphaned. Depth limits control vertical extension but place no bound on horizontal branching factor.

\medskip
\noindent\textbf{Limitation 3: Drift is a structural property, not a depth property.} Drift is defined in terms of anchor connectivity, not spawn depth. A depth limit can reduce the probability of drift occurring before human detection, but it does not eliminate the underlying condition. An agent at depth 2 with no anchor is as structurally drifted as one at depth 100. Imposing a depth limit without anchor propagation is analogous to limiting a chain's length while allowing each link to be unconnected to the next.

Formal proof: Consider any agent $a$ spawned with depth limit $D_{\max}$. Let $\anchors(a)$ denote the set of human anchors reachable from $a$ via its ancestor chain. By definition of drift: $\drifted(a) \iff \anchors(a) = \emptyset$. A depth constraint $| \anc(a) | \leq D_{\max}$ places no lower bound on $|\anchors(a)|$. Therefore, $\drifted(a)$ may hold for any $a$ with $| \anc(a) | \leq D_{\max}$. Depth limits neither prevent nor detect drift.

\subsection{Sufficient Conditions for Drift-Resistance}

The analysis above implies that drift-resistance requires structural constraints beyond depth limits:

\begin{enumerate}
  \item \textbf{Anchor propagation invariant:} Every spawned agent must receive an anchor credential in its spawn payload, regardless of depth. Formally: $\forall\, a, \child(a): \anchor(\child(a)) \neq \nil$.
  \item \textbf{Anchor liveliness check:} The anchor must be periodically verifiable. An agent whose anchor is unreachable for more than $T_{\max}$ consecutive verification cycles enters a \texttt{DRIFT\_DETECTED} state and halts.
  \item \textbf{Single-frame termination authority:} At least one anchor in every ancestor chain must have the ability to issue a \texttt{TERMINATE\_TREE} command that recursively halts all descendants, regardless of spawn depth.
\end{enumerate}

These conditions are explored in detail in Section 4, where we derive the architecture of the BX3 AgentOS framework's drift-resistance guarantees.

\subsection{Formal Model}

We consolidate the above into a compact formal model.

\medskip
\noindent\textbf{System.} A spawning system $\Sigma = (A, \rightarrow, \mathcal{A})$ where $A$ is the set of live agents, $\rightarrow$ is the spawn relation, and $\mathcal{A}: A \rightarrow \{\text{anchor\_id}\} \cup \{\nil\}$ maps each agent to its anchor (or $\nil$ if none).

\medskip
\noindent\textbf{State transition.} For any agent $a \in A$:

\begin{equation}
\state(a) = 
\begin{cases}
\active & \text{if } \mathcal{A}(a) \neq \nil \text{ and } \parent(a) \text{ is live} \\
\orphaned & \text{if } \parent(a) \text{ is } \terminated \lor \unreachable \\
\drifted & \text{if } \mathcal{A}(a) = \nil \\
\terminated & \text{if } \texttt{TERM}(a) \text{ has been executed}
\end{cases}
\end{equation}

\medskip
\noindent\textbf{Drift condition.} $\forall a \in A: \drifted(a) \iff \mathcal{A}(a) = \nil$. Equivalently:

\begin{equation}
\text{System drift} \iff \left| \{ a \in A : \mathcal{A}(a) = \nil \} \right| > 0
\end{equation}

A system with zero drifted agents is said to be \emph{anchor-saturated}. The goal of BX3 AgentOS is to maintain anchor saturation under all spawn operations and runtime conditions.

\medskip
\noindent\textbf{Theorem (Depth-limit insufficiency).} Let $D_{\max}$ be any positive integer. There exists a system $\Sigma$ such that (1) $\forall a \in A: \depth(a) \leq D_{\max}$, and (2) $\exists a \in A: \drifted(a)$. \qed

This theorem holds because depth and anchor are independent properties. A depth limit constrains tree height; it says nothing about anchor assignment. Any spawning system that fails to propagate anchors correctly will exhibit drift regardless of depth bounds.

\medskip
\noindent\textbf{Theorem (Anchor propagation sufficiency).} If $\forall a, c: (a \rightarrow c) \implies \mathcal{A}(c) \neq \nil$, then no agent in $\Sigma$ is drifted. \qed

This establishes that anchor propagation at spawn time is a sufficient condition for drift avoidance, making it the correct architectural target rather than depth limitation.